\chapter{Limites et travaux futurs}

\section{Introduction du chapitre}
Ce chapitre met en perspective les acquis du projet en identifiant les limites de l'état actuel et les évolutions prioritaires. Il permet de distinguer ce qui est valide dans le prototype de ce qui reste à consolider pour un usage de production.

\section{Limites identifiees}
Malgré les résultats obtenus, certaines limites doivent être explicitées :
\begin{itemize}
	\item déploiement orienté laboratoire, sans architecture haute disponibilité complète ;
	\item couverture de tests encore centrée sur les parcours critiques, non exhaustive sur tous les cas limites ;
	\item dépendance à des configurations manuelles pour certaines intégrations d'outils ;
	\item niveau de durcissement variable selon les environnements cibles.
\end{itemize}

\section{Pistes d'amelioration}
Les évolutions prioritaires à court terme sont :
\begin{itemize}
	\item automatiser davantage les tests de non-régression du workflow JIT ;
	\item enrichir les tableaux de bord de conformité (indicateurs périodiques, alertes) ;
	\item formaliser les procédures d'exploitation (runbooks incidents et révocation d'urgence) ;
	\item compléter la base de preuves pour les audits internes.
\end{itemize}

\section{Perspectives long terme}
A plus long terme, la solution peut évoluer vers :
\begin{itemize}
	\item une extension du périmètre à davantage d'actifs et d'équipes ;
	\item l'industrialisation complète (pipeline de déploiement, tests automatisés, gestion de secrets) ;
	\item une corrélation plus avancée entre vulnérabilités, exposition effective et identités actives ;
	\item un alignement renforcé avec les exigences GRC et les contrôles de conformité.
\end{itemize}

\section{Conclusion du chapitre}
Les limites constatées restent maîtrisables et principalement liées au niveau d'industrialisation. Les perspectives définies offrent une trajectoire claire pour faire évoluer le prototype vers une solution robuste et généralisable.
